% Options for packages loaded elsewhere
\PassOptionsToPackage{unicode}{hyperref}
\PassOptionsToPackage{hyphens}{url}
\documentclass[
]{article}
\usepackage{xcolor}
\usepackage[margin=1in]{geometry}
\usepackage{amsmath,amssymb}
\setcounter{secnumdepth}{5}
\usepackage{iftex}
\ifPDFTeX
  \usepackage[T1]{fontenc}
  \usepackage[utf8]{inputenc}
  \usepackage{textcomp} % provide euro and other symbols
\else % if luatex or xetex
  \usepackage{unicode-math} % this also loads fontspec
  \defaultfontfeatures{Scale=MatchLowercase}
  \defaultfontfeatures[\rmfamily]{Ligatures=TeX,Scale=1}
\fi
\usepackage{lmodern}
\ifPDFTeX\else
  % xetex/luatex font selection
\fi
% Use upquote if available, for straight quotes in verbatim environments
\IfFileExists{upquote.sty}{\usepackage{upquote}}{}
\IfFileExists{microtype.sty}{% use microtype if available
  \usepackage[]{microtype}
  \UseMicrotypeSet[protrusion]{basicmath} % disable protrusion for tt fonts
}{}
\makeatletter
\@ifundefined{KOMAClassName}{% if non-KOMA class
  \IfFileExists{parskip.sty}{%
    \usepackage{parskip}
  }{% else
    \setlength{\parindent}{0pt}
    \setlength{\parskip}{6pt plus 2pt minus 1pt}}
}{% if KOMA class
  \KOMAoptions{parskip=half}}
\makeatother
\usepackage{color}
\usepackage{fancyvrb}
\newcommand{\VerbBar}{|}
\newcommand{\VERB}{\Verb[commandchars=\\\{\}]}
\DefineVerbatimEnvironment{Highlighting}{Verbatim}{commandchars=\\\{\}}
% Add ',fontsize=\small' for more characters per line
\usepackage{framed}
\definecolor{shadecolor}{RGB}{248,248,248}
\newenvironment{Shaded}{\begin{snugshade}}{\end{snugshade}}
\newcommand{\AlertTok}[1]{\textcolor[rgb]{0.94,0.16,0.16}{#1}}
\newcommand{\AnnotationTok}[1]{\textcolor[rgb]{0.56,0.35,0.01}{\textbf{\textit{#1}}}}
\newcommand{\AttributeTok}[1]{\textcolor[rgb]{0.13,0.29,0.53}{#1}}
\newcommand{\BaseNTok}[1]{\textcolor[rgb]{0.00,0.00,0.81}{#1}}
\newcommand{\BuiltInTok}[1]{#1}
\newcommand{\CharTok}[1]{\textcolor[rgb]{0.31,0.60,0.02}{#1}}
\newcommand{\CommentTok}[1]{\textcolor[rgb]{0.56,0.35,0.01}{\textit{#1}}}
\newcommand{\CommentVarTok}[1]{\textcolor[rgb]{0.56,0.35,0.01}{\textbf{\textit{#1}}}}
\newcommand{\ConstantTok}[1]{\textcolor[rgb]{0.56,0.35,0.01}{#1}}
\newcommand{\ControlFlowTok}[1]{\textcolor[rgb]{0.13,0.29,0.53}{\textbf{#1}}}
\newcommand{\DataTypeTok}[1]{\textcolor[rgb]{0.13,0.29,0.53}{#1}}
\newcommand{\DecValTok}[1]{\textcolor[rgb]{0.00,0.00,0.81}{#1}}
\newcommand{\DocumentationTok}[1]{\textcolor[rgb]{0.56,0.35,0.01}{\textbf{\textit{#1}}}}
\newcommand{\ErrorTok}[1]{\textcolor[rgb]{0.64,0.00,0.00}{\textbf{#1}}}
\newcommand{\ExtensionTok}[1]{#1}
\newcommand{\FloatTok}[1]{\textcolor[rgb]{0.00,0.00,0.81}{#1}}
\newcommand{\FunctionTok}[1]{\textcolor[rgb]{0.13,0.29,0.53}{\textbf{#1}}}
\newcommand{\ImportTok}[1]{#1}
\newcommand{\InformationTok}[1]{\textcolor[rgb]{0.56,0.35,0.01}{\textbf{\textit{#1}}}}
\newcommand{\KeywordTok}[1]{\textcolor[rgb]{0.13,0.29,0.53}{\textbf{#1}}}
\newcommand{\NormalTok}[1]{#1}
\newcommand{\OperatorTok}[1]{\textcolor[rgb]{0.81,0.36,0.00}{\textbf{#1}}}
\newcommand{\OtherTok}[1]{\textcolor[rgb]{0.56,0.35,0.01}{#1}}
\newcommand{\PreprocessorTok}[1]{\textcolor[rgb]{0.56,0.35,0.01}{\textit{#1}}}
\newcommand{\RegionMarkerTok}[1]{#1}
\newcommand{\SpecialCharTok}[1]{\textcolor[rgb]{0.81,0.36,0.00}{\textbf{#1}}}
\newcommand{\SpecialStringTok}[1]{\textcolor[rgb]{0.31,0.60,0.02}{#1}}
\newcommand{\StringTok}[1]{\textcolor[rgb]{0.31,0.60,0.02}{#1}}
\newcommand{\VariableTok}[1]{\textcolor[rgb]{0.00,0.00,0.00}{#1}}
\newcommand{\VerbatimStringTok}[1]{\textcolor[rgb]{0.31,0.60,0.02}{#1}}
\newcommand{\WarningTok}[1]{\textcolor[rgb]{0.56,0.35,0.01}{\textbf{\textit{#1}}}}
\usepackage{graphicx}
\makeatletter
\newsavebox\pandoc@box
\newcommand*\pandocbounded[1]{% scales image to fit in text height/width
  \sbox\pandoc@box{#1}%
  \Gscale@div\@tempa{\textheight}{\dimexpr\ht\pandoc@box+\dp\pandoc@box\relax}%
  \Gscale@div\@tempb{\linewidth}{\wd\pandoc@box}%
  \ifdim\@tempb\p@<\@tempa\p@\let\@tempa\@tempb\fi% select the smaller of both
  \ifdim\@tempa\p@<\p@\scalebox{\@tempa}{\usebox\pandoc@box}%
  \else\usebox{\pandoc@box}%
  \fi%
}
% Set default figure placement to htbp
\def\fps@figure{htbp}
\makeatother
\setlength{\emergencystretch}{3em} % prevent overfull lines
\providecommand{\tightlist}{%
  \setlength{\itemsep}{0pt}\setlength{\parskip}{0pt}}
\usepackage{booktabs}
\usepackage{longtable}
\usepackage{array}
\usepackage{multirow}
\usepackage{wrapfig}
\usepackage{float}
\usepackage{colortbl}
\usepackage{pdflscape}
\usepackage{tabu}
\usepackage{threeparttable}
\usepackage{threeparttablex}
\usepackage[normalem]{ulem}
\usepackage{makecell}
\usepackage{xcolor}
\usepackage{bookmark}
\IfFileExists{xurl.sty}{\usepackage{xurl}}{} % add URL line breaks if available
\urlstyle{same}
\hypersetup{
  pdftitle={Causal Inference Project: G-Computation for Adversarial Multi-Agent Trust Games},
  pdfauthor={Harshit Soni (hs5666), Gabriel Nixon Raj (gr2513)},
  hidelinks,
  pdfcreator={LaTeX via pandoc}}

\title{Causal Inference Project: G-Computation for Adversarial
Multi-Agent Trust Games}
\author{Harshit Soni (hs5666), Gabriel Nixon Raj (gr2513)}
\date{2026-04-16}

\begin{document}
\maketitle

{
\setcounter{tocdepth}{3}
\tableofcontents
}
\section{Motivation}\label{motivation}

Adversarial manipulation of human decision-making is a growing concern
in domains where automated agents interact repeatedly with human users.
Dezfouli, Nock, and Dayan (2020) demonstrated that reinforcement
learning agents can discover and exploit systematic biases in human
sequential decision-making within a multi-round trust game (MRTT). Their
work raises a natural causal question: \emph{does the structure of
information sharing among adversarial agents causally affect their
ability to exploit a human decision-maker?}

We extend the Dezfouli et al.~framework by introducing a \emph{world}
variable \(W \in \{0, 1\}\) that controls whether adversarial agents can
observe each other's actions and states. In \(W = 0\) (independent),
each agent interacts with the investor through a shared learner but
cannot observe other agents' strategies. In \(W = 1\) (communication),
agents observe each other, potentially enabling coordinated
exploitation. Crucially, even under \(W = 0\), agents are not truly
independent: they are coupled through the shared investor who learns
from all agents simultaneously, creating indirect interference through a
common cause.

We chose g-computation (the parametric g-formula) as our primary
estimator for three reasons. First, our fully-specified structural
causal model (SCM) directly motivates g-computation as the natural
identification strategy under sequential ignorability. Second,
g-computation has not been covered in course homeworks, which focused on
propensity score methods, instrumental variables, and BART. Third, the
g-formula allows us to cleanly demonstrate what happens when its key
assumption --- sequential ignorability --- holds versus when it is
violated due to unobserved latent states, a distinction built directly
into our simulation design.

\section{Assumptions}\label{assumptions}

\subsection{Structural Assumptions}\label{structural-assumptions}

The SCM induces the following causal dependencies at each time step
\(t\):
\[h_t \rightarrow a_t \rightarrow x_t \rightarrow r_t \rightarrow h_{t+1}\]
with simultaneous wealth dynamics \(a_t, x_t \rightarrow w_{t+1}\) and
world-level structural modifications \(W \rightarrow h_{t,i}\) and
\(W \rightarrow x_{t,i}\). The world variable \(W\) modifies the
structural laws governing agent behavior rather than entering as an
ordinary covariate.

\subsection{Sequential Ignorability}\label{sequential-ignorability}

G-computation requires that at each time step \(t\), the treatment
assignment is ignorable given the observed history:
\[Y_t(\bar{a}_t) \perp A_t \mid X_t\] where \(X_t\) includes all
relevant state variables up to time \(t\). In our simulation, this holds
\emph{by construction} when the analyst observes the full state
including the RNN hidden states \(h_t\), since these are the sufficient
statistics for the behavioral model's policy.

\subsection{Positivity}\label{positivity}

For each level of the covariates \(X_t\), both treatment conditions must
have positive probability: \(0 < P(W = 1 \mid X_t) < 1\). Since \(W\) is
randomized by the experimenter (us) at the episode level and is
independent of all pre-treatment covariates, positivity holds trivially.

\subsection{SUTVA Considerations}\label{sutva-considerations}

The stable unit treatment value assumption is potentially violated
because agents within a game are coupled through the shared investor.
Under \(W = 0\), interference operates indirectly through the investor's
learning dynamics. Under \(W = 1\), direct agent-to-agent observation
creates additional interference channels. Our analysis operates at the
episode level where \(W\) is constant, sidestepping within-episode
interference by treating the entire episode trajectory as the outcome.

\subsection{Consistency}\label{consistency}

The consistency assumption requires that the observed outcome under
treatment \(W = w\) equals the potential outcome \(Y(w)\). Since \(W\)
is a structural parameter of the simulation --- not a choice variable
influenced by agent behavior --- consistency holds by design.

\section{Design and Estimator}\label{design-and-estimator}

\subsection{G-Computation (Parametric
G-Formula)}\label{g-computation-parametric-g-formula}

G-computation, introduced by Robins (1986), is a standardization-based
approach to causal inference under time-varying confounding. The key
idea is to model the outcome as a function of treatment and covariates,
then \emph{standardize} by computing predicted outcomes under each
treatment condition for all units.

The procedure is:

\begin{enumerate}
\def\labelenumi{\arabic{enumi}.}
\tightlist
\item
  \textbf{Fit an outcome model}: Estimate \(\hat{E}[Y \mid W, X]\) using
  linear regression with treatment--covariate interactions.
\item
  \textbf{Predict potential outcomes}: For each unit \(i\), compute
  \(\hat{Y}_i(1) = \hat{E}[Y \mid W = 1, X_i]\) and
  \(\hat{Y}_i(0) = \hat{E}[Y \mid W = 0, X_i]\).
\item
  \textbf{Estimate the ATE}:
  \(\hat{\tau} = \frac{1}{n}\sum_{i=1}^{n}[\hat{Y}_i(1) - \hat{Y}_i(0)]\).
\end{enumerate}

\subsection{Comparison with Related
Methods}\label{comparison-with-related-methods}

Unlike inverse-probability weighting (IPW), g-computation does not
require modeling the treatment assignment mechanism --- useful here
since \(W\) is randomized and the ``propensity score'' is constant at
0.5. Unlike doubly-robust estimators (AIPW), it relies entirely on
correct specification of the outcome model. This makes g-computation a
clean test case: when the outcome model is correctly specified (Setting
A), the estimator is consistent; when it is misspecified due to omitted
confounders (Setting B), we can directly observe the resulting bias.

\subsection{Implementation}\label{implementation}

We implement g-computation with linear outcome models including main
effects and treatment--covariate interactions. The covariate sets differ
across settings:

\begin{itemize}
\tightlist
\item
  \textbf{Setting A (Full Information)}: Episode number, run-level RNN
  hidden state summaries (10 dimensions), and eval-phase behavioral
  statistics. This represents the scenario where the analyst has access
  to the agents' latent internal states.
\item
  \textbf{Setting B (Behavioral Only)}: Episode number alone. This
  represents the realistic scenario (analogous to Dezfouli et al.'s
  human experiments) where the analyst observes only behavioral data and
  the agents' internal representations are latent.
\item
  \textbf{Naive Baseline}: Simple difference in means with no covariate
  adjustment.
\end{itemize}

\section{Target Estimand}\label{target-estimand}

The average treatment effect (ATE) of inter-agent communication on game
outcomes: \[\tau_{\text{ATE}} = E[Y(W = 1)] - E[Y(W = 0)]\]

We report the ATE for three outcome variables:

\begin{enumerate}
\def\labelenumi{\arabic{enumi}.}
\tightlist
\item
  \textbf{Investor cumulative return}: Total wealth accumulated by the
  investor (learner) over an episode. This captures the net effect of
  adversarial communication on the investor's welfare.
\item
  \textbf{Fair agent mean repayment fraction}: Average repayment rate by
  the fair agent. A causal effect here indicates that communication
  changes how cooperative agents behave.
\item
  \textbf{Max agent mean repayment fraction}: Average repayment rate by
  the max (exploitative) agent. A causal effect here directly measures
  whether communication enables more effective exploitation strategies.
\end{enumerate}

\section{Simulation Setup}\label{simulation-setup}

\subsection{Data Generating Process}\label{data-generating-process}

\begin{Shaded}
\begin{Highlighting}[]
\CommentTok{\# This chunk runs the Python MRTT simulation via the R driver.}
\CommentTok{\# Only executes if RUN\_DGP = TRUE in the config chunk above.}

\FunctionTok{cat}\NormalTok{(}\StringTok{"Running DGP...}\SpecialCharTok{\textbackslash{}n}\StringTok{"}\NormalTok{)}
\FunctionTok{cat}\NormalTok{(}\FunctionTok{sprintf}\NormalTok{(}\StringTok{"Seeds: \%d{-}\%d, Time budget: \%.1f hours}\SpecialCharTok{\textbackslash{}n}\StringTok{"}\NormalTok{,}
\NormalTok{            SEED\_START, SEED\_START }\SpecialCharTok{+}\NormalTok{ N\_SEEDS }\SpecialCharTok{{-}} \DecValTok{1}\NormalTok{, TIME\_BUDGET))}

\NormalTok{dgp\_cmd }\OtherTok{\textless{}{-}} \FunctionTok{sprintf}\NormalTok{(}
  \StringTok{"Rscript run\_dgp.R {-}{-}n{-}seeds \%d {-}{-}seed{-}start \%d {-}{-}time{-}budget{-}hours \%.1f {-}{-}python \%s"}\NormalTok{,}
\NormalTok{  N\_SEEDS, SEED\_START, TIME\_BUDGET, PYTHON\_CMD}
\NormalTok{)}
\FunctionTok{system}\NormalTok{(dgp\_cmd, }\AttributeTok{intern =} \ConstantTok{FALSE}\NormalTok{)}
\end{Highlighting}
\end{Shaded}

\begin{Shaded}
\begin{Highlighting}[]
\CommentTok{\# ============================================================}
\CommentTok{\# Load and stitch DGP output}
\CommentTok{\# ============================================================}

\NormalTok{dgp\_files }\OtherTok{\textless{}{-}} \FunctionTok{list.files}\NormalTok{(PROJECT\_DIR,}
                         \AttributeTok{pattern =} \StringTok{"\^{}mrtt\_dgp\_output.*}\SpecialCharTok{\textbackslash{}\textbackslash{}}\StringTok{.csv$"}\NormalTok{,}
                         \AttributeTok{full.names =} \ConstantTok{TRUE}\NormalTok{)}
\FunctionTok{stopifnot}\NormalTok{(}\StringTok{"No mrtt\_dgp\_output*.csv files found"} \OtherTok{=} \FunctionTok{length}\NormalTok{(dgp\_files) }\SpecialCharTok{\textgreater{}} \DecValTok{0}\NormalTok{)}

\NormalTok{dgp }\OtherTok{\textless{}{-}} \FunctionTok{map\_dfr}\NormalTok{(dgp\_files, }\SpecialCharTok{\textasciitilde{}} \FunctionTok{read\_csv}\NormalTok{(.x, }\AttributeTok{show\_col\_types =} \ConstantTok{FALSE}\NormalTok{)) }\SpecialCharTok{\%\textgreater{}\%}
  \FunctionTok{distinct}\NormalTok{(seed, world, episode, }\AttributeTok{.keep\_all =} \ConstantTok{TRUE}\NormalTok{) }\SpecialCharTok{\%\textgreater{}\%}
  \FunctionTok{arrange}\NormalTok{(seed, world, episode)}

\FunctionTok{cat}\NormalTok{(}\FunctionTok{sprintf}\NormalTok{(}\StringTok{"Loaded \%d rows: \%d seeds × 2 worlds × \%d episodes}\SpecialCharTok{\textbackslash{}n}\StringTok{"}\NormalTok{,}
            \FunctionTok{nrow}\NormalTok{(dgp), }\FunctionTok{n\_distinct}\NormalTok{(dgp}\SpecialCharTok{$}\NormalTok{seed),}
\NormalTok{            dgp }\SpecialCharTok{\%\textgreater{}\%} \FunctionTok{count}\NormalTok{(seed, world) }\SpecialCharTok{\%\textgreater{}\%} \FunctionTok{pull}\NormalTok{(n) }\SpecialCharTok{\%\textgreater{}\%} \FunctionTok{median}\NormalTok{()))}
\end{Highlighting}
\end{Shaded}

\begin{verbatim}
## Loaded 6800 rows: 34 seeds × 2 worlds × 100 episodes
\end{verbatim}

\subsection{Features Under
Investigation}\label{features-under-investigation}

We explore two key features of g-computation:

\begin{enumerate}
\def\labelenumi{\arabic{enumi}.}
\item
  \textbf{Reliance on sequential ignorability.} G-computation requires
  that all confounders between treatment and outcome are observed and
  conditioned upon. We test this by varying the analyst's information
  set: Setting A observes the full state (ignorability holds); Setting B
  omits the RNN hidden states (ignorability may be violated).
\item
  \textbf{Sensitivity to outcome model specification.} With randomized
  treatment, covariate adjustment should improve precision without
  introducing bias. We compare the g-computation estimates against a
  naive difference-in-means baseline to assess whether the added
  modeling complexity buys precision.
\end{enumerate}

\subsection{DGP Description}\label{dgp-description}

The simulation uses the Multi-Round Trust Task (MRTT) architecture from
Dezfouli, Nock, and Dayan (2020), extended with:

\begin{itemize}
\tightlist
\item
  \textbf{Two adversarial agents} (one fair, one max) per game
\item
  \textbf{A learning investor} (GRU-based RNN frozen from behavioral
  data, with a DQN Q-learner for investment allocation)
\item
  \textbf{10,000 rounds per episode}, 100 training episodes per run
\item
  \textbf{34 paired seeds}: each seed \(s\) generates one W0 run and one
  W1 run with identical random initialization, providing a matched-pairs
  design
\end{itemize}

\begin{Shaded}
\begin{Highlighting}[]
\NormalTok{dgp\_summary }\OtherTok{\textless{}{-}}\NormalTok{ dgp }\SpecialCharTok{\%\textgreater{}\%}
  \FunctionTok{group\_by}\NormalTok{(world) }\SpecialCharTok{\%\textgreater{}\%}
  \FunctionTok{summarise}\NormalTok{(}
    \AttributeTok{n\_episodes =} \FunctionTok{n}\NormalTok{(),}
    \AttributeTok{n\_seeds =} \FunctionTok{n\_distinct}\NormalTok{(seed),}
    \AttributeTok{mean\_inv\_return =} \FunctionTok{mean}\NormalTok{(investor\_return),}
    \AttributeTok{sd\_inv\_return =} \FunctionTok{sd}\NormalTok{(investor\_return),}
    \AttributeTok{mean\_fair\_repay =} \FunctionTok{mean}\NormalTok{(fair\_1\_repay),}
    \AttributeTok{sd\_fair\_repay =} \FunctionTok{sd}\NormalTok{(fair\_1\_repay),}
    \AttributeTok{mean\_max\_repay =} \FunctionTok{mean}\NormalTok{(max\_1\_repay),}
    \AttributeTok{sd\_max\_repay =} \FunctionTok{sd}\NormalTok{(max\_1\_repay),}
    \AttributeTok{.groups =} \StringTok{"drop"}
\NormalTok{  ) }\SpecialCharTok{\%\textgreater{}\%}
  \FunctionTok{mutate}\NormalTok{(}\AttributeTok{world\_label =} \FunctionTok{ifelse}\NormalTok{(world }\SpecialCharTok{==} \DecValTok{1}\NormalTok{,}
                               \StringTok{"W1 (Communication)"}\NormalTok{, }\StringTok{"W0 (Independent)"}\NormalTok{))}

\FunctionTok{kable}\NormalTok{(dgp\_summary }\SpecialCharTok{\%\textgreater{}\%}
        \FunctionTok{select}\NormalTok{(world\_label, n\_episodes, n\_seeds,}
\NormalTok{               mean\_inv\_return, sd\_inv\_return,}
\NormalTok{               mean\_fair\_repay, sd\_fair\_repay,}
\NormalTok{               mean\_max\_repay, sd\_max\_repay),}
      \AttributeTok{digits =} \DecValTok{2}\NormalTok{,}
      \AttributeTok{col.names =} \FunctionTok{c}\NormalTok{(}\StringTok{"World"}\NormalTok{, }\StringTok{"Episodes"}\NormalTok{, }\StringTok{"Seeds"}\NormalTok{,}
                     \StringTok{"Mean Inv. Return"}\NormalTok{, }\StringTok{"SD Inv. Return"}\NormalTok{,}
                     \StringTok{"Mean Fair Repay"}\NormalTok{, }\StringTok{"SD Fair Repay"}\NormalTok{,}
                     \StringTok{"Mean Max Repay"}\NormalTok{, }\StringTok{"SD Max Repay"}\NormalTok{),}
      \AttributeTok{caption =} \StringTok{"Summary statistics by world condition"}\NormalTok{) }\SpecialCharTok{\%\textgreater{}\%}
  \FunctionTok{kable\_styling}\NormalTok{(}\AttributeTok{bootstrap\_options =} \FunctionTok{c}\NormalTok{(}\StringTok{"striped"}\NormalTok{, }\StringTok{"hover"}\NormalTok{),}
                \AttributeTok{full\_width =} \ConstantTok{FALSE}\NormalTok{)}
\end{Highlighting}
\end{Shaded}

\begin{longtable}[t]{lrrrrrrrr}
\caption{\label{tab:dgp-summary}Summary statistics by world condition}\\
\toprule
World & Episodes & Seeds & Mean Inv. Return & SD Inv. Return & Mean Fair Repay & SD Fair Repay & Mean Max Repay & SD Max Repay\\
\midrule
W0 (Independent) & 3400 & 34 & 159203.7 & 108306.6 & 0.59 & 0.18 & 0.35 & 0.11\\
W1 (Communication) & 3400 & 34 & 155185.1 & 103748.0 & 0.59 & 0.17 & 0.35 & 0.10\\
\bottomrule
\end{longtable}

\begin{Shaded}
\begin{Highlighting}[]
\NormalTok{raw\_long }\OtherTok{\textless{}{-}}\NormalTok{ dgp }\SpecialCharTok{\%\textgreater{}\%}
  \FunctionTok{select}\NormalTok{(seed, world, episode, investor\_return, fair\_1\_repay, max\_1\_repay) }\SpecialCharTok{\%\textgreater{}\%}
  \FunctionTok{pivot\_longer}\NormalTok{(}\FunctionTok{c}\NormalTok{(investor\_return, fair\_1\_repay, max\_1\_repay),}
               \AttributeTok{names\_to =} \StringTok{"outcome"}\NormalTok{, }\AttributeTok{values\_to =} \StringTok{"value"}\NormalTok{) }\SpecialCharTok{\%\textgreater{}\%}
  \FunctionTok{mutate}\NormalTok{(}
    \AttributeTok{world\_label =} \FunctionTok{ifelse}\NormalTok{(world }\SpecialCharTok{==} \DecValTok{1}\NormalTok{, }\StringTok{"W1 (Communication)"}\NormalTok{, }\StringTok{"W0 (Independent)"}\NormalTok{),}
    \AttributeTok{outcome\_label =} \FunctionTok{case\_when}\NormalTok{(}
\NormalTok{      outcome }\SpecialCharTok{==} \StringTok{"investor\_return"} \SpecialCharTok{\textasciitilde{}} \StringTok{"Investor Cumulative Return"}\NormalTok{,}
\NormalTok{      outcome }\SpecialCharTok{==} \StringTok{"fair\_1\_repay"} \SpecialCharTok{\textasciitilde{}} \StringTok{"Fair Agent Mean Repayment"}\NormalTok{,}
\NormalTok{      outcome }\SpecialCharTok{==} \StringTok{"max\_1\_repay"} \SpecialCharTok{\textasciitilde{}} \StringTok{"Max Agent Mean Repayment"}
\NormalTok{    )}
\NormalTok{  )}

\FunctionTok{ggplot}\NormalTok{(raw\_long, }\FunctionTok{aes}\NormalTok{(}\AttributeTok{x =}\NormalTok{ value, }\AttributeTok{fill =}\NormalTok{ world\_label)) }\SpecialCharTok{+}
  \FunctionTok{geom\_histogram}\NormalTok{(}\AttributeTok{alpha =} \FloatTok{0.5}\NormalTok{, }\AttributeTok{bins =} \DecValTok{40}\NormalTok{, }\AttributeTok{position =} \StringTok{"identity"}\NormalTok{) }\SpecialCharTok{+}
  \FunctionTok{facet\_wrap}\NormalTok{(}\SpecialCharTok{\textasciitilde{}}\NormalTok{ outcome\_label, }\AttributeTok{scales =} \StringTok{"free"}\NormalTok{, }\AttributeTok{ncol =} \DecValTok{1}\NormalTok{) }\SpecialCharTok{+}
  \FunctionTok{scale\_fill\_manual}\NormalTok{(}\AttributeTok{values =} \FunctionTok{c}\NormalTok{(}\StringTok{"W0 (Independent)"} \OtherTok{=} \StringTok{"\#4575b4"}\NormalTok{,}
                                \StringTok{"W1 (Communication)"} \OtherTok{=} \StringTok{"\#d73027"}\NormalTok{)) }\SpecialCharTok{+}
  \FunctionTok{labs}\NormalTok{(}\AttributeTok{x =} \StringTok{"Value"}\NormalTok{, }\AttributeTok{y =} \StringTok{"Count"}\NormalTok{, }\AttributeTok{fill =} \StringTok{"World"}\NormalTok{) }\SpecialCharTok{+}
  \FunctionTok{theme\_minimal}\NormalTok{(}\AttributeTok{base\_size =} \DecValTok{12}\NormalTok{) }\SpecialCharTok{+}
  \FunctionTok{theme}\NormalTok{(}\AttributeTok{legend.position =} \StringTok{"bottom"}\NormalTok{, }\AttributeTok{strip.text =} \FunctionTok{element\_text}\NormalTok{(}\AttributeTok{face =} \StringTok{"bold"}\NormalTok{))}
\end{Highlighting}
\end{Shaded}

\begin{figure}

{\centering \includegraphics{causal_analysis_files/figure-latex/fig-raw-distributions-1} 

}

\caption{Outcome distributions by world condition}\label{fig:fig-raw-distributions}
\end{figure}

\subsection{Two Simulation Settings}\label{two-simulation-settings}

\begin{Shaded}
\begin{Highlighting}[]
\CommentTok{\# ============================================================}
\CommentTok{\# Extract hidden{-}state features from per{-}run game\_log.csv files}
\CommentTok{\# (Setting A covariates)}
\CommentTok{\# ============================================================}

\NormalTok{outputs\_dir }\OtherTok{\textless{}{-}} \FunctionTok{file.path}\NormalTok{(PROJECT\_DIR, }\StringTok{"outputs"}\NormalTok{)}
\NormalTok{run\_ids }\OtherTok{\textless{}{-}} \FunctionTok{unique}\NormalTok{(dgp}\SpecialCharTok{$}\NormalTok{run\_id)}

\NormalTok{parse\_hidden }\OtherTok{\textless{}{-}} \ControlFlowTok{function}\NormalTok{(hs\_str) \{}
  \ControlFlowTok{if}\NormalTok{ (}\FunctionTok{is.na}\NormalTok{(hs\_str) }\SpecialCharTok{||}\NormalTok{ hs\_str }\SpecialCharTok{==} \StringTok{""}\NormalTok{) }\FunctionTok{return}\NormalTok{(}\FunctionTok{rep}\NormalTok{(}\ConstantTok{NA\_real\_}\NormalTok{, }\DecValTok{5}\NormalTok{))}
  \FunctionTok{as.numeric}\NormalTok{(}\FunctionTok{strsplit}\NormalTok{(}\FunctionTok{gsub}\NormalTok{(}\StringTok{"}\SpecialCharTok{\textbackslash{}\textbackslash{}}\StringTok{[|}\SpecialCharTok{\textbackslash{}\textbackslash{}}\StringTok{]"}\NormalTok{, }\StringTok{""}\NormalTok{, hs\_str), }\StringTok{",}\SpecialCharTok{\textbackslash{}\textbackslash{}}\StringTok{s*"}\NormalTok{)[[}\DecValTok{1}\NormalTok{]])}
\NormalTok{\}}

\NormalTok{extract\_run\_features }\OtherTok{\textless{}{-}} \ControlFlowTok{function}\NormalTok{(rid) \{}
\NormalTok{  gl\_path }\OtherTok{\textless{}{-}} \FunctionTok{file.path}\NormalTok{(outputs\_dir, rid, }\StringTok{"game\_log.csv"}\NormalTok{)}
  \ControlFlowTok{if}\NormalTok{ (}\SpecialCharTok{!}\FunctionTok{file.exists}\NormalTok{(gl\_path)) }\FunctionTok{return}\NormalTok{(}\ConstantTok{NULL}\NormalTok{)}

  \FunctionTok{tryCatch}\NormalTok{(\{}
\NormalTok{    gl }\OtherTok{\textless{}{-}} \FunctionTok{read\_csv}\NormalTok{(gl\_path, }\AttributeTok{show\_col\_types =} \ConstantTok{FALSE}\NormalTok{)}

    \CommentTok{\# Per{-}agent behavioral summaries from eval simulation}
\NormalTok{    agent\_stats }\OtherTok{\textless{}{-}}\NormalTok{ gl }\SpecialCharTok{\%\textgreater{}\%}
      \FunctionTok{group\_by}\NormalTok{(agent\_name) }\SpecialCharTok{\%\textgreater{}\%}
      \FunctionTok{summarise}\NormalTok{(}
        \AttributeTok{mean\_repay =} \FunctionTok{mean}\NormalTok{(repayment\_pct, }\AttributeTok{na.rm =} \ConstantTok{TRUE}\NormalTok{),}
        \AttributeTok{sd\_repay =} \FunctionTok{sd}\NormalTok{(repayment\_pct, }\AttributeTok{na.rm =} \ConstantTok{TRUE}\NormalTok{),}
        \AttributeTok{mean\_invest =} \FunctionTok{mean}\NormalTok{(investment, }\AttributeTok{na.rm =} \ConstantTok{TRUE}\NormalTok{),}
        \AttributeTok{final\_hs =} \FunctionTok{last}\NormalTok{(rnn\_hidden\_state),}
        \AttributeTok{.groups =} \StringTok{"drop"}
\NormalTok{      )}

\NormalTok{    fair\_hs }\OtherTok{\textless{}{-}} \FunctionTok{parse\_hidden}\NormalTok{(}
\NormalTok{      agent\_stats }\SpecialCharTok{\%\textgreater{}\%} \FunctionTok{filter}\NormalTok{(agent\_name }\SpecialCharTok{==} \StringTok{"fair\_1"}\NormalTok{) }\SpecialCharTok{\%\textgreater{}\%} \FunctionTok{pull}\NormalTok{(final\_hs))}
\NormalTok{    max\_hs }\OtherTok{\textless{}{-}} \FunctionTok{parse\_hidden}\NormalTok{(}
\NormalTok{      agent\_stats }\SpecialCharTok{\%\textgreater{}\%} \FunctionTok{filter}\NormalTok{(agent\_name }\SpecialCharTok{==} \StringTok{"max\_1"}\NormalTok{) }\SpecialCharTok{\%\textgreater{}\%} \FunctionTok{pull}\NormalTok{(final\_hs))}

\NormalTok{    out }\OtherTok{\textless{}{-}} \FunctionTok{tibble}\NormalTok{(}
      \AttributeTok{run\_id =}\NormalTok{ rid,}
      \AttributeTok{eval\_repay\_fair =}\NormalTok{ agent\_stats }\SpecialCharTok{\%\textgreater{}\%}
        \FunctionTok{filter}\NormalTok{(agent\_name }\SpecialCharTok{==} \StringTok{"fair\_1"}\NormalTok{) }\SpecialCharTok{\%\textgreater{}\%} \FunctionTok{pull}\NormalTok{(mean\_repay),}
      \AttributeTok{eval\_repay\_max =}\NormalTok{ agent\_stats }\SpecialCharTok{\%\textgreater{}\%}
        \FunctionTok{filter}\NormalTok{(agent\_name }\SpecialCharTok{==} \StringTok{"max\_1"}\NormalTok{) }\SpecialCharTok{\%\textgreater{}\%} \FunctionTok{pull}\NormalTok{(mean\_repay),}
      \AttributeTok{eval\_sd\_fair =}\NormalTok{ agent\_stats }\SpecialCharTok{\%\textgreater{}\%}
        \FunctionTok{filter}\NormalTok{(agent\_name }\SpecialCharTok{==} \StringTok{"fair\_1"}\NormalTok{) }\SpecialCharTok{\%\textgreater{}\%} \FunctionTok{pull}\NormalTok{(sd\_repay),}
      \AttributeTok{eval\_sd\_max =}\NormalTok{ agent\_stats }\SpecialCharTok{\%\textgreater{}\%}
        \FunctionTok{filter}\NormalTok{(agent\_name }\SpecialCharTok{==} \StringTok{"max\_1"}\NormalTok{) }\SpecialCharTok{\%\textgreater{}\%} \FunctionTok{pull}\NormalTok{(sd\_repay),}
      \AttributeTok{eval\_invest =} \FunctionTok{mean}\NormalTok{(gl}\SpecialCharTok{$}\NormalTok{investment, }\AttributeTok{na.rm =} \ConstantTok{TRUE}\NormalTok{)}
\NormalTok{    )}
    \ControlFlowTok{for}\NormalTok{ (d }\ControlFlowTok{in} \FunctionTok{seq\_along}\NormalTok{(fair\_hs)) \{}
\NormalTok{      out[[}\FunctionTok{paste0}\NormalTok{(}\StringTok{"fair\_h"}\NormalTok{, d)]] }\OtherTok{\textless{}{-}}\NormalTok{ fair\_hs[d]}
\NormalTok{      out[[}\FunctionTok{paste0}\NormalTok{(}\StringTok{"max\_h"}\NormalTok{, d)]] }\OtherTok{\textless{}{-}}\NormalTok{ max\_hs[d]}
\NormalTok{    \}}
\NormalTok{    out}
\NormalTok{  \}, }\AttributeTok{error =} \ControlFlowTok{function}\NormalTok{(e) }\ConstantTok{NULL}\NormalTok{)}
\NormalTok{\}}

\NormalTok{run\_features }\OtherTok{\textless{}{-}} \FunctionTok{map\_dfr}\NormalTok{(run\_ids, extract\_run\_features)}
\NormalTok{has\_hidden }\OtherTok{\textless{}{-}} \FunctionTok{nrow}\NormalTok{(run\_features) }\SpecialCharTok{\textgreater{}} \DecValTok{0}

\ControlFlowTok{if}\NormalTok{ (has\_hidden) \{}
\NormalTok{  dgp\_full }\OtherTok{\textless{}{-}}\NormalTok{ dgp }\SpecialCharTok{\%\textgreater{}\%} \FunctionTok{left\_join}\NormalTok{(run\_features, }\AttributeTok{by =} \StringTok{"run\_id"}\NormalTok{)}
  \FunctionTok{cat}\NormalTok{(}\FunctionTok{sprintf}\NormalTok{(}\StringTok{"Hidden{-}state features extracted for \%d / \%d runs}\SpecialCharTok{\textbackslash{}n}\StringTok{"}\NormalTok{,}
              \FunctionTok{nrow}\NormalTok{(run\_features), }\FunctionTok{length}\NormalTok{(run\_ids)))}
\NormalTok{\} }\ControlFlowTok{else}\NormalTok{ \{}
\NormalTok{  dgp\_full }\OtherTok{\textless{}{-}}\NormalTok{ dgp}
  \FunctionTok{cat}\NormalTok{(}\StringTok{"No per{-}run game\_log.csv files found.}\SpecialCharTok{\textbackslash{}n}\StringTok{"}\NormalTok{)}
  \FunctionTok{cat}\NormalTok{(}\StringTok{"Setting A will use episode{-}level covariates only.}\SpecialCharTok{\textbackslash{}n}\StringTok{"}\NormalTok{)}
  \FunctionTok{cat}\NormalTok{(}\StringTok{"For full Setting A, ensure outputs/\textless{}run\_id\textgreater{}/game\_log.csv are accessible.}\SpecialCharTok{\textbackslash{}n}\StringTok{"}\NormalTok{)}
\NormalTok{\}}
\end{Highlighting}
\end{Shaded}

\begin{verbatim}
## Hidden-state features extracted for 68 / 68 runs
\end{verbatim}

\textbf{Setting A (Assumption Holds):} The analyst observes the full
state, including the RNN hidden state dimensions extracted from each
run's evaluation simulation. These hidden states capture the internal
learned representation of each agent --- the latent factors that mediate
between treatment (\(W\)) and outcome. With these observed, sequential
ignorability is satisfied by construction: there are no unobserved
common causes of treatment and outcome that aren't conditioned on.

\textbf{Setting B (Assumption Violated):} The analyst observes only the
behavioral outcomes --- episode number and the outcome variables
themselves. The RNN hidden states are treated as unobserved. This
mirrors the realistic scenario in Dezfouli et al.'s human experiments,
where the human's internal decision-making state is latent. Sequential
ignorability may fail because the hidden state influences both how the
investor allocates (related to treatment response) and the episode
outcomes.

\section{Simulation Results}\label{simulation-results}

\begin{Shaded}
\begin{Highlighting}[]
\CommentTok{\# ============================================================}
\CommentTok{\# Ground truth ATEs}
\CommentTok{\# ============================================================}

\NormalTok{outcomes }\OtherTok{\textless{}{-}} \FunctionTok{c}\NormalTok{(}\StringTok{"investor\_return"}\NormalTok{, }\StringTok{"fair\_1\_repay"}\NormalTok{, }\StringTok{"max\_1\_repay"}\NormalTok{)}
\NormalTok{outcome\_labels }\OtherTok{\textless{}{-}} \FunctionTok{c}\NormalTok{(}
  \AttributeTok{investor\_return =} \StringTok{"Investor Cumulative Return"}\NormalTok{,}
  \AttributeTok{fair\_1\_repay =} \StringTok{"Fair Agent Mean Repayment"}\NormalTok{,}
  \AttributeTok{max\_1\_repay =} \StringTok{"Max Agent Mean Repayment"}
\NormalTok{)}

\NormalTok{ground\_truth }\OtherTok{\textless{}{-}}\NormalTok{ dgp\_full }\SpecialCharTok{\%\textgreater{}\%}
  \FunctionTok{group\_by}\NormalTok{(world) }\SpecialCharTok{\%\textgreater{}\%}
  \FunctionTok{summarise}\NormalTok{(}\FunctionTok{across}\NormalTok{(}\FunctionTok{all\_of}\NormalTok{(outcomes), mean), }\AttributeTok{.groups =} \StringTok{"drop"}\NormalTok{)}

\NormalTok{ate\_true }\OtherTok{\textless{}{-}} \FunctionTok{tibble}\NormalTok{(}
  \AttributeTok{outcome =}\NormalTok{ outcomes,}
  \AttributeTok{label =}\NormalTok{ outcome\_labels[outcomes],}
  \AttributeTok{ate\_true =} \FunctionTok{map\_dbl}\NormalTok{(outcomes, }\ControlFlowTok{function}\NormalTok{(y) \{}
\NormalTok{    ground\_truth[[y]][ground\_truth}\SpecialCharTok{$}\NormalTok{world }\SpecialCharTok{==} \DecValTok{1}\NormalTok{] }\SpecialCharTok{{-}}
\NormalTok{      ground\_truth[[y]][ground\_truth}\SpecialCharTok{$}\NormalTok{world }\SpecialCharTok{==} \DecValTok{0}\NormalTok{]}
\NormalTok{  \})}
\NormalTok{)}

\FunctionTok{kable}\NormalTok{(ate\_true, }\AttributeTok{digits =} \DecValTok{4}\NormalTok{, }\AttributeTok{col.names =} \FunctionTok{c}\NormalTok{(}\StringTok{"Outcome"}\NormalTok{, }\StringTok{"Label"}\NormalTok{, }\StringTok{"True ATE"}\NormalTok{),}
      \AttributeTok{caption =} \StringTok{"Ground{-}truth ATEs from the DGP"}\NormalTok{) }\SpecialCharTok{\%\textgreater{}\%}
  \FunctionTok{kable\_styling}\NormalTok{(}\AttributeTok{bootstrap\_options =} \FunctionTok{c}\NormalTok{(}\StringTok{"striped"}\NormalTok{), }\AttributeTok{full\_width =} \ConstantTok{FALSE}\NormalTok{)}
\end{Highlighting}
\end{Shaded}

\begin{longtable}[t]{llr}
\caption{\label{tab:ground-truth}Ground-truth ATEs from the DGP}\\
\toprule
Outcome & Label & True ATE\\
\midrule
investor\_return & Investor Cumulative Return & -4018.5288\\
fair\_1\_repay & Fair Agent Mean Repayment & -0.0030\\
max\_1\_repay & Max Agent Mean Repayment & 0.0022\\
\bottomrule
\end{longtable}

\begin{Shaded}
\begin{Highlighting}[]
\CommentTok{\# ============================================================}
\CommentTok{\# G{-}Computation Estimator}
\CommentTok{\# ============================================================}

\NormalTok{gcomp\_estimate }\OtherTok{\textless{}{-}} \ControlFlowTok{function}\NormalTok{(data, outcome, covariates) \{}
\NormalTok{  rhs }\OtherTok{\textless{}{-}} \FunctionTok{paste}\NormalTok{(}\FunctionTok{c}\NormalTok{(}\StringTok{"world"}\NormalTok{, covariates), }\AttributeTok{collapse =} \StringTok{" + "}\NormalTok{)}
  \ControlFlowTok{if}\NormalTok{ (}\FunctionTok{length}\NormalTok{(covariates) }\SpecialCharTok{\textgreater{}} \DecValTok{0}\NormalTok{) \{}
\NormalTok{    interactions }\OtherTok{\textless{}{-}} \FunctionTok{paste0}\NormalTok{(}\StringTok{"world:"}\NormalTok{, covariates)}
\NormalTok{    rhs }\OtherTok{\textless{}{-}} \FunctionTok{paste}\NormalTok{(}\FunctionTok{c}\NormalTok{(rhs, interactions), }\AttributeTok{collapse =} \StringTok{" + "}\NormalTok{)}
\NormalTok{  \}}
\NormalTok{  fml }\OtherTok{\textless{}{-}} \FunctionTok{as.formula}\NormalTok{(}\FunctionTok{paste}\NormalTok{(outcome, }\StringTok{"\textasciitilde{}"}\NormalTok{, rhs))}
\NormalTok{  fit }\OtherTok{\textless{}{-}} \FunctionTok{lm}\NormalTok{(fml, }\AttributeTok{data =}\NormalTok{ data)}

\NormalTok{  d0 }\OtherTok{\textless{}{-}}\NormalTok{ d1 }\OtherTok{\textless{}{-}}\NormalTok{ data}
\NormalTok{  d0}\SpecialCharTok{$}\NormalTok{world }\OtherTok{\textless{}{-}} \DecValTok{0}
\NormalTok{  d1}\SpecialCharTok{$}\NormalTok{world }\OtherTok{\textless{}{-}} \DecValTok{1}
\NormalTok{  ate }\OtherTok{\textless{}{-}} \FunctionTok{mean}\NormalTok{(}\FunctionTok{predict}\NormalTok{(fit, }\AttributeTok{newdata =}\NormalTok{ d1) }\SpecialCharTok{{-}} \FunctionTok{predict}\NormalTok{(fit, }\AttributeTok{newdata =}\NormalTok{ d0))}
  \FunctionTok{list}\NormalTok{(}\AttributeTok{ate =}\NormalTok{ ate, }\AttributeTok{fit =}\NormalTok{ fit)}
\NormalTok{\}}

\CommentTok{\# Define covariate sets}
\NormalTok{covs\_A\_candidates }\OtherTok{\textless{}{-}} \FunctionTok{c}\NormalTok{(}\StringTok{"episode"}\NormalTok{,}
                        \StringTok{"eval\_repay\_fair"}\NormalTok{, }\StringTok{"eval\_repay\_max"}\NormalTok{,}
                        \StringTok{"eval\_sd\_fair"}\NormalTok{, }\StringTok{"eval\_sd\_max"}\NormalTok{, }\StringTok{"eval\_invest"}\NormalTok{,}
                        \FunctionTok{paste0}\NormalTok{(}\StringTok{"fair\_h"}\NormalTok{, }\DecValTok{1}\SpecialCharTok{:}\DecValTok{5}\NormalTok{), }\FunctionTok{paste0}\NormalTok{(}\StringTok{"max\_h"}\NormalTok{, }\DecValTok{1}\SpecialCharTok{:}\DecValTok{5}\NormalTok{))}

\CommentTok{\# Keep only covariates that actually exist and aren\textquotesingle{}t all{-}NA}
\NormalTok{covs\_A }\OtherTok{\textless{}{-}}\NormalTok{ covs\_A\_candidates[}
\NormalTok{  covs\_A\_candidates }\SpecialCharTok{\%in\%} \FunctionTok{names}\NormalTok{(dgp\_full) }\SpecialCharTok{\&}
    \FunctionTok{map\_lgl}\NormalTok{(covs\_A\_candidates, }\SpecialCharTok{\textasciitilde{}} \SpecialCharTok{!}\FunctionTok{all}\NormalTok{(}\FunctionTok{is.na}\NormalTok{(dgp\_full[[.x]])))}
\NormalTok{]}
\NormalTok{covs\_B }\OtherTok{\textless{}{-}} \FunctionTok{c}\NormalTok{(}\StringTok{"episode"}\NormalTok{)}

\FunctionTok{cat}\NormalTok{(}\FunctionTok{sprintf}\NormalTok{(}\StringTok{"Setting A covariates (\%d): \%s}\SpecialCharTok{\textbackslash{}n}\StringTok{"}\NormalTok{,}
            \FunctionTok{length}\NormalTok{(covs\_A), }\FunctionTok{paste}\NormalTok{(covs\_A, }\AttributeTok{collapse =} \StringTok{", "}\NormalTok{)))}
\end{Highlighting}
\end{Shaded}

\begin{verbatim}
## Setting A covariates (6): episode, eval_repay_fair, eval_repay_max, eval_sd_fair, eval_sd_max, eval_invest
\end{verbatim}

\begin{Shaded}
\begin{Highlighting}[]
\FunctionTok{cat}\NormalTok{(}\FunctionTok{sprintf}\NormalTok{(}\StringTok{"Setting B covariates (\%d): \%s}\SpecialCharTok{\textbackslash{}n}\StringTok{"}\NormalTok{,}
            \FunctionTok{length}\NormalTok{(covs\_B), }\FunctionTok{paste}\NormalTok{(covs\_B, }\AttributeTok{collapse =} \StringTok{", "}\NormalTok{)))}
\end{Highlighting}
\end{Shaded}

\begin{verbatim}
## Setting B covariates (1): episode
\end{verbatim}

\begin{Shaded}
\begin{Highlighting}[]
\CommentTok{\# ============================================================}
\CommentTok{\# Point estimates (full sample)}
\CommentTok{\# ============================================================}

\NormalTok{point\_ests }\OtherTok{\textless{}{-}} \FunctionTok{map\_dfr}\NormalTok{(outcomes, }\ControlFlowTok{function}\NormalTok{(y) \{}
\NormalTok{  w1 }\OtherTok{\textless{}{-}} \FunctionTok{mean}\NormalTok{(dgp\_full[[y]][dgp\_full}\SpecialCharTok{$}\NormalTok{world }\SpecialCharTok{==} \DecValTok{1}\NormalTok{], }\AttributeTok{na.rm =} \ConstantTok{TRUE}\NormalTok{)}
\NormalTok{  w0 }\OtherTok{\textless{}{-}} \FunctionTok{mean}\NormalTok{(dgp\_full[[y]][dgp\_full}\SpecialCharTok{$}\NormalTok{world }\SpecialCharTok{==} \DecValTok{0}\NormalTok{], }\AttributeTok{na.rm =} \ConstantTok{TRUE}\NormalTok{)}
  \FunctionTok{tibble}\NormalTok{(}
    \AttributeTok{outcome =}\NormalTok{ y,}
    \AttributeTok{naive   =}\NormalTok{ w1 }\SpecialCharTok{{-}}\NormalTok{ w0,}
    \AttributeTok{gcomp\_A =} \FunctionTok{gcomp\_estimate}\NormalTok{(dgp\_full, y, covs\_A)}\SpecialCharTok{$}\NormalTok{ate,}
    \AttributeTok{gcomp\_B =} \FunctionTok{gcomp\_estimate}\NormalTok{(dgp\_full, y, covs\_B)}\SpecialCharTok{$}\NormalTok{ate}
\NormalTok{  )}
\NormalTok{\})}

\FunctionTok{kable}\NormalTok{(point\_ests, }\AttributeTok{digits =} \DecValTok{4}\NormalTok{,}
      \AttributeTok{col.names =} \FunctionTok{c}\NormalTok{(}\StringTok{"Outcome"}\NormalTok{, }\StringTok{"Naive DiM"}\NormalTok{, }\StringTok{"G{-}comp A (Full)"}\NormalTok{, }\StringTok{"G{-}comp B (Behav.)"}\NormalTok{),}
      \AttributeTok{caption =} \StringTok{"Point estimates from full sample"}\NormalTok{) }\SpecialCharTok{\%\textgreater{}\%}
  \FunctionTok{kable\_styling}\NormalTok{(}\AttributeTok{bootstrap\_options =} \FunctionTok{c}\NormalTok{(}\StringTok{"striped"}\NormalTok{), }\AttributeTok{full\_width =} \ConstantTok{FALSE}\NormalTok{)}
\end{Highlighting}
\end{Shaded}

\begin{longtable}[t]{lrrr}
\caption{\label{tab:point-estimates}Point estimates from full sample}\\
\toprule
Outcome & Naive DiM & G-comp A (Full) & G-comp B (Behav.)\\
\midrule
investor\_return & -4018.5288 & -1660.3995 & -4018.5288\\
fair\_1\_repay & -0.0030 & -0.0029 & -0.0030\\
max\_1\_repay & 0.0022 & 0.0036 & 0.0022\\
\bottomrule
\end{longtable}

\begin{Shaded}
\begin{Highlighting}[]
\CommentTok{\# ============================================================}
\CommentTok{\# Cluster Bootstrap (seed{-}level)}
\CommentTok{\# ============================================================}

\NormalTok{all\_seeds }\OtherTok{\textless{}{-}} \FunctionTok{sort}\NormalTok{(}\FunctionTok{unique}\NormalTok{(dgp\_full}\SpecialCharTok{$}\NormalTok{seed))}
\NormalTok{n\_seeds }\OtherTok{\textless{}{-}} \FunctionTok{length}\NormalTok{(all\_seeds)}

\NormalTok{boot\_one }\OtherTok{\textless{}{-}} \ControlFlowTok{function}\NormalTok{(b) \{}
\NormalTok{  boot\_seeds }\OtherTok{\textless{}{-}} \FunctionTok{sample}\NormalTok{(all\_seeds, n\_seeds, }\AttributeTok{replace =} \ConstantTok{TRUE}\NormalTok{)}
\NormalTok{  boot\_data }\OtherTok{\textless{}{-}} \FunctionTok{map2\_dfr}\NormalTok{(boot\_seeds, }\FunctionTok{seq\_along}\NormalTok{(boot\_seeds), }\ControlFlowTok{function}\NormalTok{(s, idx) \{}
\NormalTok{    dgp\_full }\SpecialCharTok{\%\textgreater{}\%} \FunctionTok{filter}\NormalTok{(seed }\SpecialCharTok{==}\NormalTok{ s) }\SpecialCharTok{\%\textgreater{}\%} \FunctionTok{mutate}\NormalTok{(}\AttributeTok{boot\_id =}\NormalTok{ idx)}
\NormalTok{  \})}

  \FunctionTok{map\_dfr}\NormalTok{(outcomes, }\ControlFlowTok{function}\NormalTok{(y) \{}
\NormalTok{    w1 }\OtherTok{\textless{}{-}} \FunctionTok{mean}\NormalTok{(boot\_data[[y]][boot\_data}\SpecialCharTok{$}\NormalTok{world }\SpecialCharTok{==} \DecValTok{1}\NormalTok{], }\AttributeTok{na.rm =} \ConstantTok{TRUE}\NormalTok{)}
\NormalTok{    w0 }\OtherTok{\textless{}{-}} \FunctionTok{mean}\NormalTok{(boot\_data[[y]][boot\_data}\SpecialCharTok{$}\NormalTok{world }\SpecialCharTok{==} \DecValTok{0}\NormalTok{], }\AttributeTok{na.rm =} \ConstantTok{TRUE}\NormalTok{)}
\NormalTok{    ate\_naive }\OtherTok{\textless{}{-}}\NormalTok{ w1 }\SpecialCharTok{{-}}\NormalTok{ w0}

\NormalTok{    ate\_A }\OtherTok{\textless{}{-}} \FunctionTok{tryCatch}\NormalTok{(}
      \FunctionTok{gcomp\_estimate}\NormalTok{(boot\_data, y, covs\_A)}\SpecialCharTok{$}\NormalTok{ate,}
      \AttributeTok{error =} \ControlFlowTok{function}\NormalTok{(e) }\ConstantTok{NA\_real\_}
\NormalTok{    )}
\NormalTok{    ate\_B }\OtherTok{\textless{}{-}} \FunctionTok{tryCatch}\NormalTok{(}
      \FunctionTok{gcomp\_estimate}\NormalTok{(boot\_data, y, covs\_B)}\SpecialCharTok{$}\NormalTok{ate,}
      \AttributeTok{error =} \ControlFlowTok{function}\NormalTok{(e) }\ConstantTok{NA\_real\_}
\NormalTok{    )}

    \FunctionTok{tibble}\NormalTok{(}\AttributeTok{boot =}\NormalTok{ b, }\AttributeTok{outcome =}\NormalTok{ y,}
           \AttributeTok{ate\_naive =}\NormalTok{ ate\_naive, }\AttributeTok{ate\_A =}\NormalTok{ ate\_A, }\AttributeTok{ate\_B =}\NormalTok{ ate\_B)}
\NormalTok{  \})}
\NormalTok{\}}

\FunctionTok{cat}\NormalTok{(}\FunctionTok{sprintf}\NormalTok{(}\StringTok{"Running \%d bootstrap replicates (seed{-}level clustering)...}\SpecialCharTok{\textbackslash{}n}\StringTok{"}\NormalTok{, N\_BOOT))}
\end{Highlighting}
\end{Shaded}

\begin{verbatim}
## Running 1000 bootstrap replicates (seed-level clustering)...
\end{verbatim}

\begin{Shaded}
\begin{Highlighting}[]
\NormalTok{t0 }\OtherTok{\textless{}{-}} \FunctionTok{Sys.time}\NormalTok{()}
\NormalTok{boot\_results }\OtherTok{\textless{}{-}} \FunctionTok{map\_dfr}\NormalTok{(}\FunctionTok{seq\_len}\NormalTok{(N\_BOOT), }\ControlFlowTok{function}\NormalTok{(b) \{}
  \ControlFlowTok{if}\NormalTok{ (b }\SpecialCharTok{\%\%} \DecValTok{200} \SpecialCharTok{==} \DecValTok{0}\NormalTok{) }\FunctionTok{cat}\NormalTok{(}\FunctionTok{sprintf}\NormalTok{(}\StringTok{"  [\%d / \%d]}\SpecialCharTok{\textbackslash{}n}\StringTok{"}\NormalTok{, b, N\_BOOT))}
  \FunctionTok{boot\_one}\NormalTok{(b)}
\NormalTok{\})}
\end{Highlighting}
\end{Shaded}

\begin{verbatim}
##   [200 / 1000]
##   [400 / 1000]
##   [600 / 1000]
##   [800 / 1000]
##   [1000 / 1000]
\end{verbatim}

\begin{Shaded}
\begin{Highlighting}[]
\NormalTok{elapsed }\OtherTok{\textless{}{-}} \FunctionTok{difftime}\NormalTok{(}\FunctionTok{Sys.time}\NormalTok{(), t0, }\AttributeTok{units =} \StringTok{"mins"}\NormalTok{)}
\FunctionTok{cat}\NormalTok{(}\FunctionTok{sprintf}\NormalTok{(}\StringTok{"Bootstrap complete in \%.1f minutes.}\SpecialCharTok{\textbackslash{}n}\StringTok{"}\NormalTok{, }\FunctionTok{as.numeric}\NormalTok{(elapsed)))}
\end{Highlighting}
\end{Shaded}

\begin{verbatim}
## Bootstrap complete in 1.2 minutes.
\end{verbatim}

\begin{Shaded}
\begin{Highlighting}[]
\CommentTok{\# ============================================================}
\CommentTok{\# Evaluation: Bias, RMSE, Coverage}
\CommentTok{\# ============================================================}

\NormalTok{estimators }\OtherTok{\textless{}{-}} \FunctionTok{c}\NormalTok{(}\StringTok{"ate\_naive"}\NormalTok{, }\StringTok{"ate\_A"}\NormalTok{, }\StringTok{"ate\_B"}\NormalTok{)}
\NormalTok{est\_labels }\OtherTok{\textless{}{-}} \FunctionTok{c}\NormalTok{(}\AttributeTok{ate\_naive =} \StringTok{"Naive DiM"}\NormalTok{,}
                \AttributeTok{ate\_A =} \StringTok{"G{-}comp (Full Info)"}\NormalTok{,}
                \AttributeTok{ate\_B =} \StringTok{"G{-}comp (Behavioral Only)"}\NormalTok{)}

\NormalTok{eval\_estimator }\OtherTok{\textless{}{-}} \ControlFlowTok{function}\NormalTok{(ates, truth, }\AttributeTok{alpha =} \FloatTok{0.05}\NormalTok{) \{}
\NormalTok{  ates }\OtherTok{\textless{}{-}}\NormalTok{ ates[}\SpecialCharTok{!}\FunctionTok{is.na}\NormalTok{(ates)]}
  \ControlFlowTok{if}\NormalTok{ (}\FunctionTok{length}\NormalTok{(ates) }\SpecialCharTok{==} \DecValTok{0}\NormalTok{) }\FunctionTok{return}\NormalTok{(}\FunctionTok{tibble}\NormalTok{(}
    \AttributeTok{bias =} \ConstantTok{NA}\NormalTok{, }\AttributeTok{rmse =} \ConstantTok{NA}\NormalTok{, }\AttributeTok{coverage =} \ConstantTok{NA}\NormalTok{,}
    \AttributeTok{mean\_est =} \ConstantTok{NA}\NormalTok{, }\AttributeTok{se =} \ConstantTok{NA}\NormalTok{, }\AttributeTok{ci\_lo =} \ConstantTok{NA}\NormalTok{, }\AttributeTok{ci\_hi =} \ConstantTok{NA}\NormalTok{))}

\NormalTok{  mn   }\OtherTok{\textless{}{-}} \FunctionTok{mean}\NormalTok{(ates)}
\NormalTok{  se   }\OtherTok{\textless{}{-}} \FunctionTok{sd}\NormalTok{(ates)}
\NormalTok{  bias }\OtherTok{\textless{}{-}}\NormalTok{ mn }\SpecialCharTok{{-}}\NormalTok{ truth}
\NormalTok{  rmse }\OtherTok{\textless{}{-}} \FunctionTok{sqrt}\NormalTok{(}\FunctionTok{mean}\NormalTok{((ates }\SpecialCharTok{{-}}\NormalTok{ truth)}\SpecialCharTok{\^{}}\DecValTok{2}\NormalTok{))}
\NormalTok{  z    }\OtherTok{\textless{}{-}} \FunctionTok{qnorm}\NormalTok{(}\DecValTok{1} \SpecialCharTok{{-}}\NormalTok{ alpha }\SpecialCharTok{/} \DecValTok{2}\NormalTok{)}

  \CommentTok{\# Coverage: fraction of bootstrap CIs (ate\_b ± z*se) containing truth}
\NormalTok{  coverage }\OtherTok{\textless{}{-}} \FunctionTok{mean}\NormalTok{((ates }\SpecialCharTok{{-}}\NormalTok{ z }\SpecialCharTok{*}\NormalTok{ se }\SpecialCharTok{\textless{}=}\NormalTok{ truth) }\SpecialCharTok{\&}\NormalTok{ (ates }\SpecialCharTok{+}\NormalTok{ z }\SpecialCharTok{*}\NormalTok{ se }\SpecialCharTok{\textgreater{}=}\NormalTok{ truth))}

  \FunctionTok{tibble}\NormalTok{(}\AttributeTok{bias =}\NormalTok{ bias, }\AttributeTok{rmse =}\NormalTok{ rmse, }\AttributeTok{coverage =}\NormalTok{ coverage,}
         \AttributeTok{mean\_est =}\NormalTok{ mn, }\AttributeTok{se =}\NormalTok{ se,}
         \AttributeTok{ci\_lo =} \FunctionTok{quantile}\NormalTok{(ates, alpha }\SpecialCharTok{/} \DecValTok{2}\NormalTok{),}
         \AttributeTok{ci\_hi =} \FunctionTok{quantile}\NormalTok{(ates, }\DecValTok{1} \SpecialCharTok{{-}}\NormalTok{ alpha }\SpecialCharTok{/} \DecValTok{2}\NormalTok{))}
\NormalTok{\}}

\NormalTok{eval\_table }\OtherTok{\textless{}{-}} \FunctionTok{map\_dfr}\NormalTok{(outcomes, }\ControlFlowTok{function}\NormalTok{(y) \{}
\NormalTok{  truth }\OtherTok{\textless{}{-}}\NormalTok{ ate\_true}\SpecialCharTok{$}\NormalTok{ate\_true[ate\_true}\SpecialCharTok{$}\NormalTok{outcome }\SpecialCharTok{==}\NormalTok{ y]}
  \FunctionTok{map\_dfr}\NormalTok{(estimators, }\ControlFlowTok{function}\NormalTok{(est) \{}
\NormalTok{    ates }\OtherTok{\textless{}{-}}\NormalTok{ boot\_results }\SpecialCharTok{\%\textgreater{}\%} \FunctionTok{filter}\NormalTok{(outcome }\SpecialCharTok{==}\NormalTok{ y) }\SpecialCharTok{\%\textgreater{}\%} \FunctionTok{pull}\NormalTok{(}\SpecialCharTok{!!}\FunctionTok{sym}\NormalTok{(est))}
\NormalTok{    metrics }\OtherTok{\textless{}{-}} \FunctionTok{eval\_estimator}\NormalTok{(ates, truth, ALPHA)}
    \FunctionTok{bind\_cols}\NormalTok{(}\FunctionTok{tibble}\NormalTok{(}\AttributeTok{outcome =}\NormalTok{ y, }\AttributeTok{estimator =}\NormalTok{ est, }\AttributeTok{truth =}\NormalTok{ truth), metrics)}
\NormalTok{  \})}
\NormalTok{\}) }\SpecialCharTok{\%\textgreater{}\%}
  \FunctionTok{mutate}\NormalTok{(}
    \AttributeTok{outcome\_label =}\NormalTok{ outcome\_labels[outcome],}
    \AttributeTok{estimator\_label =}\NormalTok{ est\_labels[estimator]}
\NormalTok{  )}
\end{Highlighting}
\end{Shaded}

\subsection{Evaluation Table}\label{evaluation-table}

\begin{Shaded}
\begin{Highlighting}[]
\NormalTok{eval\_display }\OtherTok{\textless{}{-}}\NormalTok{ eval\_table }\SpecialCharTok{\%\textgreater{}\%}
  \FunctionTok{select}\NormalTok{(outcome\_label, estimator\_label, truth, mean\_est,}
\NormalTok{         bias, rmse, coverage, se) }\SpecialCharTok{\%\textgreater{}\%}
  \FunctionTok{arrange}\NormalTok{(outcome\_label, estimator\_label)}

\FunctionTok{kable}\NormalTok{(eval\_display, }\AttributeTok{digits =} \DecValTok{4}\NormalTok{,}
      \AttributeTok{col.names =} \FunctionTok{c}\NormalTok{(}\StringTok{"Outcome"}\NormalTok{, }\StringTok{"Estimator"}\NormalTok{, }\StringTok{"True ATE"}\NormalTok{, }\StringTok{"Mean Est."}\NormalTok{,}
                     \StringTok{"Bias"}\NormalTok{, }\StringTok{"RMSE"}\NormalTok{, }\StringTok{"Coverage"}\NormalTok{, }\StringTok{"SE"}\NormalTok{),}
      \AttributeTok{caption =} \StringTok{"Bias, RMSE, and 95}\SpecialCharTok{\textbackslash{}\textbackslash{}}\StringTok{\% CI Coverage across estimators and outcomes"}\NormalTok{) }\SpecialCharTok{\%\textgreater{}\%}
  \FunctionTok{kable\_styling}\NormalTok{(}\AttributeTok{bootstrap\_options =} \FunctionTok{c}\NormalTok{(}\StringTok{"striped"}\NormalTok{, }\StringTok{"hover"}\NormalTok{),}
                \AttributeTok{full\_width =} \ConstantTok{FALSE}\NormalTok{, }\AttributeTok{font\_size =} \DecValTok{11}\NormalTok{) }\SpecialCharTok{\%\textgreater{}\%}
  \FunctionTok{column\_spec}\NormalTok{(}\DecValTok{5}\NormalTok{, }\AttributeTok{bold =} \ConstantTok{TRUE}\NormalTok{)}
\end{Highlighting}
\end{Shaded}

\begingroup\fontsize{11}{13}\selectfont

\begin{longtable}[t]{llrr>{}rrrr}
\caption{\label{tab:eval-table-display}Bias, RMSE, and 95\% CI Coverage across estimators and outcomes}\\
\toprule
Outcome & Estimator & True ATE & Mean Est. & Bias & RMSE & Coverage & SE\\
\midrule
Fair Agent Mean Repayment & G-comp (Behavioral Only) & -0.0030 & -0.0030 & \textbf{0.0000} & 0.0045 & 0.953 & 0.0045\\
Fair Agent Mean Repayment & G-comp (Full Info) & -0.0030 & -0.0025 & \textbf{0.0005} & 0.0049 & 0.946 & 0.0049\\
Fair Agent Mean Repayment & Naive DiM & -0.0030 & -0.0030 & \textbf{0.0000} & 0.0045 & 0.953 & 0.0045\\
Investor Cumulative Return & G-comp (Behavioral Only) & -4018.5288 & -4125.7334 & \textbf{-107.2046} & 3176.9729 & 0.954 & 3176.7524\\
Investor Cumulative Return & G-comp (Full Info) & -4018.5288 & -1474.3972 & \textbf{2544.1316} & 4531.2748 & 0.894 & 3751.5223\\
\addlinespace
Investor Cumulative Return & Naive DiM & -4018.5288 & -4125.7334 & \textbf{-107.2046} & 3176.9729 & 0.954 & 3176.7524\\
Max Agent Mean Repayment & G-comp (Behavioral Only) & 0.0022 & 0.0020 & \textbf{-0.0001} & 0.0046 & 0.951 & 0.0046\\
Max Agent Mean Repayment & G-comp (Full Info) & 0.0022 & 0.0033 & \textbf{0.0011} & 0.0067 & 0.944 & 0.0066\\
Max Agent Mean Repayment & Naive DiM & 0.0022 & 0.0020 & \textbf{-0.0001} & 0.0046 & 0.951 & 0.0046\\
\bottomrule
\end{longtable}
\endgroup{}

\subsection{Bootstrap Distributions}\label{bootstrap-distributions}

\begin{Shaded}
\begin{Highlighting}[]
\NormalTok{boot\_long }\OtherTok{\textless{}{-}}\NormalTok{ boot\_results }\SpecialCharTok{\%\textgreater{}\%}
  \FunctionTok{pivot\_longer}\NormalTok{(}\FunctionTok{starts\_with}\NormalTok{(}\StringTok{"ate\_"}\NormalTok{),}
               \AttributeTok{names\_to =} \StringTok{"estimator"}\NormalTok{, }\AttributeTok{values\_to =} \StringTok{"ate\_hat"}\NormalTok{) }\SpecialCharTok{\%\textgreater{}\%}
  \FunctionTok{mutate}\NormalTok{(}
    \AttributeTok{estimator\_label =}\NormalTok{ est\_labels[estimator],}
    \AttributeTok{outcome\_label =}\NormalTok{ outcome\_labels[outcome]}
\NormalTok{  )}

\FunctionTok{ggplot}\NormalTok{(boot\_long, }\FunctionTok{aes}\NormalTok{(}\AttributeTok{x =}\NormalTok{ ate\_hat, }\AttributeTok{fill =}\NormalTok{ estimator\_label)) }\SpecialCharTok{+}
  \FunctionTok{geom\_histogram}\NormalTok{(}\AttributeTok{alpha =} \FloatTok{0.6}\NormalTok{, }\AttributeTok{bins =} \DecValTok{50}\NormalTok{, }\AttributeTok{position =} \StringTok{"identity"}\NormalTok{) }\SpecialCharTok{+}
  \FunctionTok{geom\_vline}\NormalTok{(}
    \AttributeTok{data =}\NormalTok{ ate\_true }\SpecialCharTok{\%\textgreater{}\%} \FunctionTok{mutate}\NormalTok{(}\AttributeTok{outcome\_label =}\NormalTok{ outcome\_labels[outcome]),}
    \FunctionTok{aes}\NormalTok{(}\AttributeTok{xintercept =}\NormalTok{ ate\_true),}
    \AttributeTok{color =} \StringTok{"red"}\NormalTok{, }\AttributeTok{linetype =} \StringTok{"dashed"}\NormalTok{, }\AttributeTok{linewidth =} \DecValTok{1}
\NormalTok{  ) }\SpecialCharTok{+}
  \FunctionTok{facet\_wrap}\NormalTok{(}\SpecialCharTok{\textasciitilde{}}\NormalTok{ outcome\_label, }\AttributeTok{scales =} \StringTok{"free"}\NormalTok{, }\AttributeTok{ncol =} \DecValTok{1}\NormalTok{) }\SpecialCharTok{+}
  \FunctionTok{scale\_fill\_brewer}\NormalTok{(}\AttributeTok{palette =} \StringTok{"Set2"}\NormalTok{) }\SpecialCharTok{+}
  \FunctionTok{labs}\NormalTok{(}\AttributeTok{x =} \StringTok{"Estimated ATE"}\NormalTok{, }\AttributeTok{y =} \StringTok{"Count"}\NormalTok{, }\AttributeTok{fill =} \StringTok{"Estimator"}\NormalTok{) }\SpecialCharTok{+}
  \FunctionTok{theme\_minimal}\NormalTok{(}\AttributeTok{base\_size =} \DecValTok{12}\NormalTok{) }\SpecialCharTok{+}
  \FunctionTok{theme}\NormalTok{(}\AttributeTok{legend.position =} \StringTok{"bottom"}\NormalTok{, }\AttributeTok{strip.text =} \FunctionTok{element\_text}\NormalTok{(}\AttributeTok{face =} \StringTok{"bold"}\NormalTok{))}
\end{Highlighting}
\end{Shaded}

\begin{figure}

{\centering \includegraphics{causal_analysis_files/figure-latex/fig-bootstrap-1} 

}

\caption{Bootstrap distributions of ATE estimates. Red dashed line marks the ground-truth ATE.}\label{fig:fig-bootstrap}
\end{figure}

\subsection{Bias Comparison}\label{bias-comparison}

\begin{Shaded}
\begin{Highlighting}[]
\FunctionTok{ggplot}\NormalTok{(eval\_table, }\FunctionTok{aes}\NormalTok{(}\AttributeTok{x =}\NormalTok{ estimator\_label, }\AttributeTok{y =}\NormalTok{ bias, }\AttributeTok{fill =}\NormalTok{ estimator\_label)) }\SpecialCharTok{+}
  \FunctionTok{geom\_col}\NormalTok{(}\AttributeTok{alpha =} \FloatTok{0.8}\NormalTok{, }\AttributeTok{width =} \FloatTok{0.6}\NormalTok{) }\SpecialCharTok{+}
  \FunctionTok{geom\_hline}\NormalTok{(}\AttributeTok{yintercept =} \DecValTok{0}\NormalTok{, }\AttributeTok{linetype =} \StringTok{"dashed"}\NormalTok{, }\AttributeTok{color =} \StringTok{"gray40"}\NormalTok{) }\SpecialCharTok{+}
  \FunctionTok{facet\_wrap}\NormalTok{(}\SpecialCharTok{\textasciitilde{}}\NormalTok{ outcome\_label, }\AttributeTok{scales =} \StringTok{"free\_y"}\NormalTok{, }\AttributeTok{ncol =} \DecValTok{1}\NormalTok{) }\SpecialCharTok{+}
  \FunctionTok{scale\_fill\_brewer}\NormalTok{(}\AttributeTok{palette =} \StringTok{"Set2"}\NormalTok{) }\SpecialCharTok{+}
  \FunctionTok{labs}\NormalTok{(}\AttributeTok{x =} \ConstantTok{NULL}\NormalTok{, }\AttributeTok{y =} \StringTok{"Bias (mean estimate {-} true ATE)"}\NormalTok{, }\AttributeTok{fill =} \StringTok{"Estimator"}\NormalTok{) }\SpecialCharTok{+}
  \FunctionTok{theme\_minimal}\NormalTok{(}\AttributeTok{base\_size =} \DecValTok{12}\NormalTok{) }\SpecialCharTok{+}
  \FunctionTok{theme}\NormalTok{(}\AttributeTok{legend.position =} \StringTok{"bottom"}\NormalTok{,}
        \AttributeTok{strip.text =} \FunctionTok{element\_text}\NormalTok{(}\AttributeTok{face =} \StringTok{"bold"}\NormalTok{),}
        \AttributeTok{axis.text.x =} \FunctionTok{element\_text}\NormalTok{(}\AttributeTok{angle =} \DecValTok{15}\NormalTok{, }\AttributeTok{hjust =} \DecValTok{1}\NormalTok{))}
\end{Highlighting}
\end{Shaded}

\begin{figure}

{\centering \includegraphics{causal_analysis_files/figure-latex/fig-bias-1} 

}

\caption{Bias by estimator and outcome. Dashed line at zero.}\label{fig:fig-bias}
\end{figure}

\subsection{RMSE Comparison}\label{rmse-comparison}

\begin{Shaded}
\begin{Highlighting}[]
\FunctionTok{ggplot}\NormalTok{(eval\_table, }\FunctionTok{aes}\NormalTok{(}\AttributeTok{x =}\NormalTok{ estimator\_label, }\AttributeTok{y =}\NormalTok{ rmse, }\AttributeTok{fill =}\NormalTok{ estimator\_label)) }\SpecialCharTok{+}
  \FunctionTok{geom\_col}\NormalTok{(}\AttributeTok{alpha =} \FloatTok{0.8}\NormalTok{, }\AttributeTok{width =} \FloatTok{0.6}\NormalTok{) }\SpecialCharTok{+}
  \FunctionTok{facet\_wrap}\NormalTok{(}\SpecialCharTok{\textasciitilde{}}\NormalTok{ outcome\_label, }\AttributeTok{scales =} \StringTok{"free\_y"}\NormalTok{, }\AttributeTok{ncol =} \DecValTok{1}\NormalTok{) }\SpecialCharTok{+}
  \FunctionTok{scale\_fill\_brewer}\NormalTok{(}\AttributeTok{palette =} \StringTok{"Set2"}\NormalTok{) }\SpecialCharTok{+}
  \FunctionTok{labs}\NormalTok{(}\AttributeTok{x =} \ConstantTok{NULL}\NormalTok{, }\AttributeTok{y =} \StringTok{"Root Mean Squared Error"}\NormalTok{, }\AttributeTok{fill =} \StringTok{"Estimator"}\NormalTok{) }\SpecialCharTok{+}
  \FunctionTok{theme\_minimal}\NormalTok{(}\AttributeTok{base\_size =} \DecValTok{12}\NormalTok{) }\SpecialCharTok{+}
  \FunctionTok{theme}\NormalTok{(}\AttributeTok{legend.position =} \StringTok{"bottom"}\NormalTok{,}
        \AttributeTok{strip.text =} \FunctionTok{element\_text}\NormalTok{(}\AttributeTok{face =} \StringTok{"bold"}\NormalTok{),}
        \AttributeTok{axis.text.x =} \FunctionTok{element\_text}\NormalTok{(}\AttributeTok{angle =} \DecValTok{15}\NormalTok{, }\AttributeTok{hjust =} \DecValTok{1}\NormalTok{))}
\end{Highlighting}
\end{Shaded}

\begin{figure}

{\centering \includegraphics{causal_analysis_files/figure-latex/fig-rmse-1} 

}

\caption{RMSE by estimator and outcome.}\label{fig:fig-rmse}
\end{figure}

\subsection{Coverage}\label{coverage}

\begin{Shaded}
\begin{Highlighting}[]
\FunctionTok{ggplot}\NormalTok{(eval\_table, }\FunctionTok{aes}\NormalTok{(}\AttributeTok{x =}\NormalTok{ estimator\_label, }\AttributeTok{y =}\NormalTok{ coverage, }\AttributeTok{fill =}\NormalTok{ estimator\_label)) }\SpecialCharTok{+}
  \FunctionTok{geom\_col}\NormalTok{(}\AttributeTok{alpha =} \FloatTok{0.8}\NormalTok{, }\AttributeTok{width =} \FloatTok{0.6}\NormalTok{) }\SpecialCharTok{+}
  \FunctionTok{geom\_hline}\NormalTok{(}\AttributeTok{yintercept =} \DecValTok{1} \SpecialCharTok{{-}}\NormalTok{ ALPHA, }\AttributeTok{linetype =} \StringTok{"dashed"}\NormalTok{, }\AttributeTok{color =} \StringTok{"red"}\NormalTok{) }\SpecialCharTok{+}
  \FunctionTok{facet\_wrap}\NormalTok{(}\SpecialCharTok{\textasciitilde{}}\NormalTok{ outcome\_label, }\AttributeTok{ncol =} \DecValTok{1}\NormalTok{) }\SpecialCharTok{+}
  \FunctionTok{scale\_fill\_brewer}\NormalTok{(}\AttributeTok{palette =} \StringTok{"Set2"}\NormalTok{) }\SpecialCharTok{+}
  \FunctionTok{scale\_y\_continuous}\NormalTok{(}\AttributeTok{limits =} \FunctionTok{c}\NormalTok{(}\DecValTok{0}\NormalTok{, }\DecValTok{1}\NormalTok{)) }\SpecialCharTok{+}
  \FunctionTok{labs}\NormalTok{(}\AttributeTok{x =} \ConstantTok{NULL}\NormalTok{, }\AttributeTok{y =} \StringTok{"Coverage Rate"}\NormalTok{, }\AttributeTok{fill =} \StringTok{"Estimator"}\NormalTok{) }\SpecialCharTok{+}
  \FunctionTok{theme\_minimal}\NormalTok{(}\AttributeTok{base\_size =} \DecValTok{12}\NormalTok{) }\SpecialCharTok{+}
  \FunctionTok{theme}\NormalTok{(}\AttributeTok{legend.position =} \StringTok{"bottom"}\NormalTok{,}
        \AttributeTok{strip.text =} \FunctionTok{element\_text}\NormalTok{(}\AttributeTok{face =} \StringTok{"bold"}\NormalTok{),}
        \AttributeTok{axis.text.x =} \FunctionTok{element\_text}\NormalTok{(}\AttributeTok{angle =} \DecValTok{15}\NormalTok{, }\AttributeTok{hjust =} \DecValTok{1}\NormalTok{))}
\end{Highlighting}
\end{Shaded}

\begin{figure}

{\centering \includegraphics{causal_analysis_files/figure-latex/fig-coverage-1} 

}

\caption{95\% CI coverage. Red dashed line marks the nominal 95\% level.}\label{fig:fig-coverage}
\end{figure}

\subsection{ATE by Training Episode}\label{ate-by-training-episode}

\begin{Shaded}
\begin{Highlighting}[]
\NormalTok{episode\_means }\OtherTok{\textless{}{-}}\NormalTok{ dgp\_full }\SpecialCharTok{\%\textgreater{}\%}
  \FunctionTok{group\_by}\NormalTok{(episode, world) }\SpecialCharTok{\%\textgreater{}\%}
  \FunctionTok{summarise}\NormalTok{(}\FunctionTok{across}\NormalTok{(}\FunctionTok{all\_of}\NormalTok{(outcomes), mean), }\AttributeTok{.groups =} \StringTok{"drop"}\NormalTok{)}

\NormalTok{ep\_w0 }\OtherTok{\textless{}{-}}\NormalTok{ episode\_means }\SpecialCharTok{\%\textgreater{}\%} \FunctionTok{filter}\NormalTok{(world }\SpecialCharTok{==} \DecValTok{0}\NormalTok{)}
\NormalTok{ep\_w1 }\OtherTok{\textless{}{-}}\NormalTok{ episode\_means }\SpecialCharTok{\%\textgreater{}\%} \FunctionTok{filter}\NormalTok{(world }\SpecialCharTok{==} \DecValTok{1}\NormalTok{)}

\NormalTok{episode\_ates }\OtherTok{\textless{}{-}} \FunctionTok{tibble}\NormalTok{(}
  \AttributeTok{episode =}\NormalTok{ ep\_w0}\SpecialCharTok{$}\NormalTok{episode,}
  \AttributeTok{ate\_investor =}\NormalTok{ ep\_w1}\SpecialCharTok{$}\NormalTok{investor\_return }\SpecialCharTok{{-}}\NormalTok{ ep\_w0}\SpecialCharTok{$}\NormalTok{investor\_return,}
  \AttributeTok{ate\_fair =}\NormalTok{ ep\_w1}\SpecialCharTok{$}\NormalTok{fair\_1\_repay }\SpecialCharTok{{-}}\NormalTok{ ep\_w0}\SpecialCharTok{$}\NormalTok{fair\_1\_repay,}
  \AttributeTok{ate\_max =}\NormalTok{ ep\_w1}\SpecialCharTok{$}\NormalTok{max\_1\_repay }\SpecialCharTok{{-}}\NormalTok{ ep\_w0}\SpecialCharTok{$}\NormalTok{max\_1\_repay}
\NormalTok{) }\SpecialCharTok{\%\textgreater{}\%}
  \FunctionTok{pivot\_longer}\NormalTok{(}\SpecialCharTok{{-}}\NormalTok{episode, }\AttributeTok{names\_to =} \StringTok{"outcome"}\NormalTok{, }\AttributeTok{values\_to =} \StringTok{"ate"}\NormalTok{) }\SpecialCharTok{\%\textgreater{}\%}
  \FunctionTok{mutate}\NormalTok{(}\AttributeTok{outcome\_label =} \FunctionTok{case\_when}\NormalTok{(}
\NormalTok{    outcome }\SpecialCharTok{==} \StringTok{"ate\_investor"} \SpecialCharTok{\textasciitilde{}} \StringTok{"Investor Cumulative Return"}\NormalTok{,}
\NormalTok{    outcome }\SpecialCharTok{==} \StringTok{"ate\_fair"} \SpecialCharTok{\textasciitilde{}} \StringTok{"Fair Agent Mean Repayment"}\NormalTok{,}
\NormalTok{    outcome }\SpecialCharTok{==} \StringTok{"ate\_max"} \SpecialCharTok{\textasciitilde{}} \StringTok{"Max Agent Mean Repayment"}
\NormalTok{  ))}

\FunctionTok{ggplot}\NormalTok{(episode\_ates, }\FunctionTok{aes}\NormalTok{(}\AttributeTok{x =}\NormalTok{ episode, }\AttributeTok{y =}\NormalTok{ ate)) }\SpecialCharTok{+}
  \FunctionTok{geom\_line}\NormalTok{(}\AttributeTok{alpha =} \FloatTok{0.3}\NormalTok{) }\SpecialCharTok{+}
  \FunctionTok{geom\_smooth}\NormalTok{(}\AttributeTok{method =} \StringTok{"loess"}\NormalTok{, }\AttributeTok{se =} \ConstantTok{TRUE}\NormalTok{, }\AttributeTok{color =} \StringTok{"steelblue"}\NormalTok{, }\AttributeTok{linewidth =} \FloatTok{1.2}\NormalTok{) }\SpecialCharTok{+}
  \FunctionTok{geom\_hline}\NormalTok{(}\AttributeTok{yintercept =} \DecValTok{0}\NormalTok{, }\AttributeTok{linetype =} \StringTok{"dashed"}\NormalTok{, }\AttributeTok{color =} \StringTok{"gray40"}\NormalTok{) }\SpecialCharTok{+}
  \FunctionTok{facet\_wrap}\NormalTok{(}\SpecialCharTok{\textasciitilde{}}\NormalTok{ outcome\_label, }\AttributeTok{scales =} \StringTok{"free\_y"}\NormalTok{, }\AttributeTok{ncol =} \DecValTok{1}\NormalTok{) }\SpecialCharTok{+}
  \FunctionTok{labs}\NormalTok{(}\AttributeTok{x =} \StringTok{"Training Episode"}\NormalTok{, }\AttributeTok{y =} \StringTok{"ATE (W1 {-} W0)"}\NormalTok{) }\SpecialCharTok{+}
  \FunctionTok{theme\_minimal}\NormalTok{(}\AttributeTok{base\_size =} \DecValTok{12}\NormalTok{) }\SpecialCharTok{+}
  \FunctionTok{theme}\NormalTok{(}\AttributeTok{strip.text =} \FunctionTok{element\_text}\NormalTok{(}\AttributeTok{face =} \StringTok{"bold"}\NormalTok{))}
\end{Highlighting}
\end{Shaded}

\begin{figure}

{\centering \includegraphics{causal_analysis_files/figure-latex/fig-ate-episode-1} 

}

\caption{ATE by training episode, showing how the causal effect of communication evolves over the course of learning.}\label{fig:fig-ate-episode}
\end{figure}

\subsection{Paired Seed Analysis}\label{paired-seed-analysis}

\begin{Shaded}
\begin{Highlighting}[]
\NormalTok{seed\_by\_world }\OtherTok{\textless{}{-}}\NormalTok{ dgp\_full }\SpecialCharTok{\%\textgreater{}\%}
  \FunctionTok{group\_by}\NormalTok{(seed, world) }\SpecialCharTok{\%\textgreater{}\%}
  \FunctionTok{summarise}\NormalTok{(}\FunctionTok{across}\NormalTok{(}\FunctionTok{all\_of}\NormalTok{(outcomes), mean), }\AttributeTok{.groups =} \StringTok{"drop"}\NormalTok{)}

\NormalTok{sw0 }\OtherTok{\textless{}{-}}\NormalTok{ seed\_by\_world }\SpecialCharTok{\%\textgreater{}\%} \FunctionTok{filter}\NormalTok{(world }\SpecialCharTok{==} \DecValTok{0}\NormalTok{) }\SpecialCharTok{\%\textgreater{}\%} \FunctionTok{arrange}\NormalTok{(seed)}
\NormalTok{sw1 }\OtherTok{\textless{}{-}}\NormalTok{ seed\_by\_world }\SpecialCharTok{\%\textgreater{}\%} \FunctionTok{filter}\NormalTok{(world }\SpecialCharTok{==} \DecValTok{1}\NormalTok{) }\SpecialCharTok{\%\textgreater{}\%} \FunctionTok{arrange}\NormalTok{(seed)}

\NormalTok{seed\_means }\OtherTok{\textless{}{-}} \FunctionTok{tibble}\NormalTok{(}
  \AttributeTok{seed =}\NormalTok{ sw0}\SpecialCharTok{$}\NormalTok{seed,}
  \AttributeTok{diff\_investor =}\NormalTok{ sw1}\SpecialCharTok{$}\NormalTok{investor\_return }\SpecialCharTok{{-}}\NormalTok{ sw0}\SpecialCharTok{$}\NormalTok{investor\_return,}
  \AttributeTok{diff\_fair =}\NormalTok{ sw1}\SpecialCharTok{$}\NormalTok{fair\_1\_repay }\SpecialCharTok{{-}}\NormalTok{ sw0}\SpecialCharTok{$}\NormalTok{fair\_1\_repay,}
  \AttributeTok{diff\_max =}\NormalTok{ sw1}\SpecialCharTok{$}\NormalTok{max\_1\_repay }\SpecialCharTok{{-}}\NormalTok{ sw0}\SpecialCharTok{$}\NormalTok{max\_1\_repay}
\NormalTok{)}

\NormalTok{seed\_long }\OtherTok{\textless{}{-}}\NormalTok{ seed\_means }\SpecialCharTok{\%\textgreater{}\%}
  \FunctionTok{pivot\_longer}\NormalTok{(}\SpecialCharTok{{-}}\NormalTok{seed, }\AttributeTok{names\_to =} \StringTok{"outcome"}\NormalTok{, }\AttributeTok{values\_to =} \StringTok{"diff"}\NormalTok{) }\SpecialCharTok{\%\textgreater{}\%}
  \FunctionTok{mutate}\NormalTok{(}\AttributeTok{outcome\_label =} \FunctionTok{case\_when}\NormalTok{(}
\NormalTok{    outcome }\SpecialCharTok{==} \StringTok{"diff\_investor"} \SpecialCharTok{\textasciitilde{}} \StringTok{"Investor Cumulative Return"}\NormalTok{,}
\NormalTok{    outcome }\SpecialCharTok{==} \StringTok{"diff\_fair"} \SpecialCharTok{\textasciitilde{}} \StringTok{"Fair Agent Mean Repayment"}\NormalTok{,}
\NormalTok{    outcome }\SpecialCharTok{==} \StringTok{"diff\_max"} \SpecialCharTok{\textasciitilde{}} \StringTok{"Max Agent Mean Repayment"}
\NormalTok{  ))}

\FunctionTok{ggplot}\NormalTok{(seed\_long, }\FunctionTok{aes}\NormalTok{(}\AttributeTok{x =} \FunctionTok{factor}\NormalTok{(seed), }\AttributeTok{y =}\NormalTok{ diff)) }\SpecialCharTok{+}
  \FunctionTok{geom\_col}\NormalTok{(}\AttributeTok{fill =} \StringTok{"steelblue"}\NormalTok{, }\AttributeTok{alpha =} \FloatTok{0.7}\NormalTok{) }\SpecialCharTok{+}
  \FunctionTok{geom\_hline}\NormalTok{(}\AttributeTok{yintercept =} \DecValTok{0}\NormalTok{, }\AttributeTok{linetype =} \StringTok{"dashed"}\NormalTok{, }\AttributeTok{color =} \StringTok{"gray40"}\NormalTok{) }\SpecialCharTok{+}
  \FunctionTok{facet\_wrap}\NormalTok{(}\SpecialCharTok{\textasciitilde{}}\NormalTok{ outcome\_label, }\AttributeTok{scales =} \StringTok{"free\_y"}\NormalTok{, }\AttributeTok{ncol =} \DecValTok{1}\NormalTok{) }\SpecialCharTok{+}
  \FunctionTok{labs}\NormalTok{(}\AttributeTok{x =} \StringTok{"Seed"}\NormalTok{, }\AttributeTok{y =} \StringTok{"Within{-}Seed Difference (W1 {-} W0)"}\NormalTok{) }\SpecialCharTok{+}
  \FunctionTok{theme\_minimal}\NormalTok{(}\AttributeTok{base\_size =} \DecValTok{11}\NormalTok{) }\SpecialCharTok{+}
  \FunctionTok{theme}\NormalTok{(}\AttributeTok{strip.text =} \FunctionTok{element\_text}\NormalTok{(}\AttributeTok{face =} \StringTok{"bold"}\NormalTok{),}
        \AttributeTok{axis.text.x =} \FunctionTok{element\_text}\NormalTok{(}\AttributeTok{size =} \DecValTok{7}\NormalTok{))}
\end{Highlighting}
\end{Shaded}

\begin{figure}

{\centering \includegraphics{causal_analysis_files/figure-latex/fig-paired-1} 

}

\caption{Paired within-seed ATEs across all seeds. Each point is one seed; bar shows the paired difference.}\label{fig:fig-paired}
\end{figure}

\section{Discussion}\label{discussion}

The simulation results reveal several patterns relevant to the choice of
causal inference methodology in multi-agent settings.

\textbf{Effect size.} The ground-truth ATE of inter-agent communication
on investor returns is modest (-4018.5), consistent with the hypothesis
that in the two-agent case, the primary interference channel operates
indirectly through the shared investor rather than through direct
agent-to-agent coordination. The paired seed analysis shows substantial
heterogeneity: 15 seeds show positive effects and 19 show negative
effects, suggesting that the direction of the communication effect
depends on the specific learning trajectory realized under each seed.

\textbf{Estimator performance.} {[}Interpret based on actual results
after knitting.{]}

\textbf{Setting A vs.~Setting B.} {[}Compare bias, RMSE, and coverage
between the full-information and behavioral-only g-computation
settings.{]}

\textbf{Non-stationarity.} The ATE-by-episode figure illustrates how the
causal effect evolves over the course of learning. Early episodes, where
agents are still exploring (\(\epsilon\) is high), may show different
treatment effects than late episodes where policies have converged. This
non-stationarity is a genuine feature of the DGP --- not an artifact ---
and highlights why sequential methods like g-computation are
appropriate: they can condition on episode number as a proxy for
learning stage.

\textbf{Limitations.} The Setting A covariates are run-level summaries
extracted from the evaluation simulation, not true per-episode hidden
states from the training loop. This means Setting A captures
\emph{between-run} variation in what the RNN learned, but not
\emph{within-run} episode-to-episode hidden state dynamics. A fully
realized Setting A would require logging \(h_t\) at each training
episode, which is a straightforward engineering extension but was not
implemented for this project. Consequently, the Setting A vs.~Setting B
contrast may understate the true benefit of observing hidden states.

\textbf{Implications.} In adversarial multi-agent environments, the
choice of causal estimator should be guided by what the analyst can
observe. When internal agent states are accessible (as in simulation),
g-computation with full conditioning provides unbiased estimates. When
only behavioral traces are available (as in real human experiments), the
analyst faces a bias--variance tradeoff that depends on the strength of
the latent confounding.

\section{References}\label{references}

\begin{enumerate}
\def\labelenumi{\arabic{enumi}.}
\item
  Dezfouli, A., Nock, R., \& Dayan, P. (2020). Adversarial
  vulnerabilities of human decision-making. \emph{PNAS}, 117(46),
  29221--29228.
\item
  Robins, J. M. (1986). A new approach to causal inference in mortality
  studies with a sustained exposure period. \emph{Mathematical
  Modelling}, 7(9-12), 1393--1512.
\item
  Bongers, S., Forré, P., Peters, J., \& Mooij, J. M. (2021).
  Foundations of structural causal models with cycles and latent
  variables. \emph{Annals of Statistics}, 49(5), 2885--2915.
\item
  Nobandegani, A. S., Rish, I., \& Shultz, T. R. (2023). Towards
  machines that trust: AI agents learn to trust in the trust game.
  \emph{arXiv:2312.12868}.
\item
  Park, S., Lu, W., \& Yang, S. (2025). Evaluating and learning optimal
  dynamic treatment regimes under truncation by death. \emph{NeurIPS
  2025}.
\end{enumerate}

\end{document}
